% Agentic Coding for Economists — Padova five-hour workshop
% Build: latexmk -pdf workshop_slides.tex  (from slides/workshop/)

\input{../shared/beamer_preamble}
\newcommand{\modebox}[1]{\fcolorbox{courseblue}{white}{\small\textsf{#1}}}

\renewcommand{\WorkshopRepoUrl}{\url{https://github.com/meleantonio/AC4E_Padova}}
\renewcommand{\ParentCourseUrl}{\url{https://github.com/meleantonio/AC4E_Pavia}}
\title{Agentic Coding for Economists}
\subtitle{Padova Workshop — Research Article Harness}
\author{Antonio Mele}
\institute{University of Padova}
\date{25 June 2026}

\begin{document}

% =============================================================================
\section{Title}
% =============================================================================

\begin{frame}{}
  \begin{center}
    {\LARGE \textbf{Agentic Coding for Economists}} \\
    \vspace{0.5em}
    {\large Five-hour hands-on workshop} \\
    \vspace{0.5em}
    {\small Participant repo: \WorkshopRepoUrl} \\
    \vspace{0.75em}
    Antonio Mele \\
    University of Padova \\
    \vspace{0.5em}
  \end{center}
\end{frame}

\begin{frame}{How We Teach Today}
  \begin{block}{Demo-first rhythm}
    \begin{enumerate}
      \item \textbf{Short demo ($\leq$8 min):} \textnormal{One tactical workflow, live.}
      \item \textbf{Hands-on sprint:} \textnormal{You produce a file in \emph{your} \artifact{examples/starter\_article/} copy.}
      \item \textbf{Debrief:} \textnormal{What did the agent do vs what did you sign off?}
    \end{enumerate}
  \end{block}
  \vspace{0.3em}
  % \begin{itemize}
  %   \item \textbf{One mini repo all day:} \textnormal{Not scattered exercises in the course repo only.}
  %   \item \textbf{Economics thread:} \textnormal{Remote-work DiD demo; identification is teachable, not outsourced.}
  %   \item \textbf{Ethics:} \textnormal{No confidential referee manuscripts in AI tools.}
  % \end{itemize}
\end{frame}

% =============================================================================
\section{Agentic Coding}
% =============================================================================

\begin{frame}{What is Agentic Coding?}
  \begin{block}{Definition}
    \concept{Agentic coding} = delegating implementation work to AI agents that \textbf{read}, \textbf{plan}, \textbf{edit}, and \textbf{verify}---while you orchestrate scope, context, and quality gates.
  \end{block}
  \vspace{0.4em}
  \begin{itemize}
    \item \textbf{Agents} \textnormal{operate on your codebase with \emph{tools}: edit files, run terminal commands, search the repo (and optionally the web).}
    \item \textbf{You} \textnormal{control what they see (context files, attachments, web resources) and how they behave (the harness: \artifact{AGENTS.md}, rules, skills, hooks, ...).}
    \item \textbf{Output} \textnormal{must be \emph{verifiable}: tests pass, replication is valid, acceptance criteria achieved.}
  \end{itemize}
\end{frame}

\begin{frame}{The Agent Loop}
  \begin{center}
    \begin{tikzpicture}[
        node distance=1.1cm,
        loopbox/.style={rectangle, draw, rounded corners, minimum width=2.1cm, minimum height=0.65cm, align=center, font=\small},
        arr/.style={-Stealth, thick}
      ]
      \node[loopbox, fill=courseblue!15] (read) {Read\\context};
      \node[loopbox, right=of read] (plan) {Plan};
      \node[loopbox, right=of plan] (act) {Act\\tools};
      \node[loopbox, right=of act] (verify) {Verify};
      \draw[arr] (read) -- (plan) -- (act) -- (verify);
      \draw[arr] (verify.north) .. controls +(0,0.55) and +(0,0.55) .. (read.north);
    \end{tikzpicture}
  \end{center}
  \vspace{0.35em}
  {\small
  \begin{itemize}
    \item \textbf{Read} \textnormal{--- Project brief, \artifact{AGENTS.md}, attached files, search results. Bad context $\Rightarrow$ bad code/output.}
    \item \textbf{Plan} \textnormal{--- Break work into steps (Plan mode, specifications, requirements, tasks, explicit goal).}
    \item \textbf{Act} \textnormal{--- Edit code, run code/tests, execute terminal commands, update paper write up.}
    \item \textbf{Verify} \textnormal{--- Re-run scripts, check replication pipeline, etc.}
  \end{itemize}}
\end{frame}

\begin{frame}{Why Economists Need This}
  \begin{itemize}
    \item \textbf{Pipeline complexity:} \textnormal{Data $\to$ cleaning $\to$ estimation $\to$ tables $\to$ paper---dozens of small, automatable steps.}
    \item \textbf{Reproducibility pressure:} \textnormal{Journals and replicators expect documented data, portable paths, one-command runs.}
    \item \textbf{Multi-artifact projects:} \textnormal{Code, LaTeX, appendix, and replication package must stay consistent.}
    \item \textbf{This workshop:} \textnormal{Build a \emph{harness} you can fork onto a research project, not just one-off chat prompts.}
  \end{itemize}
\end{frame}

\begin{frame}{What is a Research Article Harness?}
  \begin{block}{Harness}
    A \textbf{project template} where context files, skills, subagents, hooks, and verification gates work together across the research lifecycle, from idea to replication.
  \end{block}
  \vspace{0.3em}
  
  \vspace{0.3em}
  \begin{itemize}
    \item \textbf{Not} \textnormal{a single magic prompt.}
    \item \textbf{Is} \textnormal{repeatable infrastructure: the agent always knows the repo map, boundaries, and how to run checks.}
    \item \textbf{Can} \textnormal{be adapted to other research projects.}
  \end{itemize}
\end{frame}

\begin{frame}{Harness Lifecycle (Workshop Spine)}
  \begin{center}
    \begin{tikzpicture}[
        node distance=0.55cm and 0.35cm,
        phase/.style={rectangle, draw, rounded corners, minimum height=0.75cm, align=center, font=\scriptsize},
        arr/.style={-Stealth, thick}
      ]
      \node[phase, fill=courseblue!15, minimum width=2.1cm] (s) {Scaffold\\Module 1};
      \node[phase, fill=courseteal!15, minimum width=2.1cm, right=of s] (h) {Harness\\Module 2};
      \node[phase, fill=courseblue!15, minimum width=2.1cm, right=of h] (p) {Pipeline\\Module 3};
      \node[phase, fill=courseteal!15, minimum width=2.1cm, right=of p] (q) {Quality\\Module 4};
      \node[phase, fill=courseblue!15, minimum width=2.1cm, right=of q] (o) {Sign-off\\Module 5};
      \draw[arr] (s) -- (h) -- (p) -- (q) -- (o);
    \end{tikzpicture}
  \end{center}
  \vspace{0.35em}
  {\small
  \begin{itemize}
    \item \textbf{Scaffold} \textnormal{--- Create repo, project idea/brief, \artifact{AGENTS.md}, privacy boundaries, three tasks.}
    \item \textbf{Harness} \textnormal{--- Create skills, subagents, hooks, MCPs; run one long goal with stop conditions.}
    \item \textbf{Pipeline} \textnormal{--- Update idea/brief, write and maintain code and paper.}
    \item \textbf{Quality} \textnormal{--- Referee report, replication check}
    \item \textbf{Sign-off} \textnormal{--- Checklist + 7-day adoption plan for your own research.}
  \end{itemize}}
\end{frame}

\begin{frame}[fragile]{Demo: Coding Agent Interface}

  
      \textbf{A typical coding agent UI:}
      \begin{itemize}
        \item \textbf{Prompt/Chat:} Type a goal\\ (e.g., ``Clean \texttt{data.csv} and compute summary stats.'')
        \item \textbf{Context} loaded: repo map, \texttt{AGENTS.md}, \artifact{templates/project\_brief.md}.
        \item \textbf{Artifacts:} View/edit code, instructions, logs, or generated files.
        \item \textbf{Tools:} Run code, preview plots, test outputs {\scriptsize(e.g., buttons for ``Run All'', ``Test'', ``View PDF'').}
      \end{itemize}
      
  
  
\end{frame}

% =============================================================================
\section{Harness vocabulary}
% =============================================================================

\begin{frame}{Project Context: \artifact{AGENTS.md}}
  \begin{block}{What it is}
    The ``front door'' for any agent: a human-readable brief at the repo root that loads every session.
  \end{block}
  \textbf{Typically includes:}
  \begin{itemize}
    \item \textbf{Research question} \textnormal{and stage (seminar draft, JMP, etc.)}
    \item \textbf{Repository map} \textnormal{--- which folders agents may read or write}
    \item \textbf{Verification commands} \textnormal{--- \texttt{pytest}, \texttt{src/estimate\_did.py}, \texttt{latexmk}}
    \item \textbf{Do-not rules} \textnormal{--- e.g.\ never edit \artifact{data/raw/} without permission}
  \end{itemize}
  \vspace{0.25em}
  \textbf{Claude lane:} \textnormal{may also use \artifact{CLAUDE.md}; keep \artifact{AGENTS.md} for collaborators.}
\end{frame}

\begin{frame}[fragile,t,shrink=5]{Example: \artifact{AGENTS.md} (HANK Macro)}
  {\footnotesize Same pattern for applied micro, structural, or macro repos.}
  \vspace{0.2em}
  \begin{lstlisting}[basicstyle=\ttfamily, frame=single, backgroundcolor=\color{coursebg}]
# AGENTS.md - HANK macro (illustrative)

## Project
HANK: monetary policy with heterogeneous households; seminar draft.

## Repo map
Read: src/, models/, data/processed/, docs/
Write: src/, models/, docs/drafts/, output/
Do NOT: data/raw/, output/confidential/, secrets/

## Verify
pytest src/ | python models/run_HANK.py | latexmk -pdf docs/paper.tex

## Rules
No raw-data edits without permission; no API keys or microdata in context.
Tests must pass; document model changes in docs/; use a branch for review.
  \end{lstlisting}
\end{frame}


\begin{frame}{Exercise: Create Your Own \artifact{AGENTS.md}}
  \textbf{Exercise:} Draft an \artifact{AGENTS.md} for your own project.
  \vspace{1em}

  \begin{block}{Include the following sections}
    \begin{itemize}
      \item \textbf{Project:} Briefly state your project's research topic and stage.
      \item \textbf{Repo map:} \textit{Which folders should your agent read or write?} What should be off-limits?
      \item \textbf{Verify:} Specify at least one command to check code or outputs (e.g., \texttt{pytest}, \texttt{make}).
      \item \textbf{Rules:} Any do-not rules, data boundaries, or requirements for contributors/agents?
    \end{itemize}
  \end{block}

  \vspace{.5em}
  {\small
    \textbf{Tip:} Imagine onboarding a new research assistant (human or AI)---what context and boundaries do they need on their first day?}

  \vspace{1em}
  \textbf{(Start from \artifact{templates/AGENTS\_economics.md} or the HANK macro example above.)}
\end{frame}


\begin{frame}{Research Spec: Brief and Design Memo}
  \begin{columns}[T,onlytextwidth]
    \begin{column}{0.48\textwidth}
      \textbf{Project Brief: \\ \artifact{templates/project\_brief.md}}
      \begin{itemize}
        \item \textbf{What:} \textnormal{Scope, audience, in/out of scope}
        \item \textbf{Why:} \textnormal{Stops agents from ``helpfully'' expanding the paper}
        \item \textbf{When:} \textnormal{At the beginning of the project, Module 1 --- personalize or keep demo}
      \end{itemize}
      {\scriptsize Example: \artifact{examples/starter\_article/docs/project\_brief.md}}
    \end{column}
    \begin{column}{0.48\textwidth}
      \textbf{Design Memo: \\ \artifact{templates/research\_design\_memo.md}}
      \begin{itemize}
        \item \textbf{What:} \textnormal{Hypotheses, data, identification, risks/threats}
        \item \textbf{Why:} \textnormal{Plan before code; hub for SDD-lite}
        \item \textbf{When:} \textnormal{Beginning of the project, but maintained throughout, Module 3 --- identification section is human-owned}
      \end{itemize}
      {\scriptsize Example: \artifact{examples/starter\_article/docs/research\_design\_memo.md}}
    \end{column}
  \end{columns}
  \vspace{1.35em}
  \textbf{Spec-Driven Development: \\ \artifact{examples/starter\_article/spec/intent.md}} \textnormal{--- optional; for ambitious projects (intent $\to$ requirements $\to$ design $\to$ tasks).}
\end{frame}

\begin{frame}{Exercise: Create Your Own Project Brief and Design Memo}
  \textbf{Exercise:} Draft a project brief and design memo for your own project.
  \vspace{0.5em}
  {\small Start from \artifact{templates/project\_brief.md} and \artifact{templates/research\_design\_memo.md}.}
  \vspace{0.5em}
  \begin{block}{Project Brief}
    \begin{itemize}
      \item \textbf{What:} \textnormal{Scope, idea,audience, in/out of scope}
      
    \end{itemize}
  \end{block}

  \begin{block}{Design Memo}
    \begin{itemize}
      \item \textbf{What:} \textnormal{Hypotheses, data, identification, risks/threats} \end{itemize}
  \end{block}
  
\end{frame}


\begin{frame}[fragile]{Privacy Boundary: \artifact{.cursorignore} or \artifact{AGENTS.md}}
  \begin{columns}[T,onlytextwidth]
    \begin{column}{0.48\textwidth}
      \textbf{What it does:} \textnormal{Controls what the agent \emph{indexes} and may send as context.}
      \begin{itemize}
        \item \textbf{Privacy:} \textnormal{API keys, microdata, confidential drafts}
        \item \textbf{Performance:} \textnormal{Large CSVs, PDFs, \texttt{.venv}}
      \end{itemize}
      \vspace{0.25em}
      \textbf{Different from \texttt{.gitignore}:} \\
      \texttt{.gitignore} = version control; \\
      \artifact{.cursorignore} = AI context.
    \end{column}
    \begin{column}{0.48\textwidth}
      \begin{lstlisting}[style=bash, basicstyle=\ttfamily\tiny]
data/raw/confidential/
.env
*.dta
*_confidential*
      \end{lstlisting}
    \end{column}
  \end{columns}
\end{frame}

\begin{frame}{What Are Agent Skills?}
  \textbf{Skills} extend the agent with \textbf{reusable workflows} stored as \artifact{SKILL.md} files.
  \begin{itemize}
    \item \textbf{Trigger:} \textnormal{The skill's \texttt{description} in YAML frontmatter matches your request.}
    \item \textbf{vs subagents:} \textnormal{Skills augment the \emph{main} agent---no separate context window.}
    \item \textbf{vs rules:} \textnormal{Rules are always-on constraints; skills are procedures you invoke.}
  \end{itemize}
  
\end{frame}

\begin{frame}{Exercise: Create Your Own Agent Skills}
  \textbf{Exercise:} Create an agent skill for your own project.
  \vspace{1em}
  \begin{block}{Skill}
    \begin{itemize}
      \item \textbf{What:} \textnormal{Description, trigger, steps}
    \end{itemize}
  \end{block}
\end{frame}

\begin{frame}{What Are Subagents?}
  \textbf{Subagents} are specialized assistants with their own \textbf{context window} and prompt.
  \begin{itemize}
    \item \textbf{Why:} \textnormal{Context isolation---a reviewer should not share the same ``implementation bias'' as the coder.}
    \item \textbf{Format:} \textnormal{\artifact{examples/starter\_article/.cursor/agents/*.md} or \artifact{examples/claude/.claude/agents/*.md} with YAML frontmatter + role prompt.}
    \item \textbf{Invoke:} \textnormal{``Use identification-reviewer on the design memo only'' or tool-specific \texttt{/agents}.}
  \end{itemize}
  \vspace{0.35em}
  
\end{frame}


\begin{frame}{Exercise: Create Your Own Subagent}
  \textbf{Exercise:} Create a subagent for your own project.
  \vspace{1em}
  \begin{block}{Subagent}
    \begin{itemize}
      \item \textbf{What:} \textnormal{Description, trigger, steps}
    \end{itemize}
  \end{block}
\end{frame}
\begin{frame}{What Are Hooks?}
  \textbf{Hooks} run \textbf{deterministic} scripts at lifecycle events (session start, before/after tool use).
  \begin{itemize}
    \item \textbf{Why:} \textnormal{LLMs forget; hooks enforce reminders (privacy, no \texttt{rm -rf}, run tests after edits).}
    \item \textbf{Where:} \textnormal{\artifact{examples/hooks/.codex/}, Claude \texttt{settings.json}, \artifact{examples/starter\_article/.cursor/hooks.json}.}
    \item \textbf{Trust:} \textnormal{Codex \texttt{/hooks}; Claude settings UI---you approve what runs.}
  \end{itemize}
  \vspace{0.3em}
  \textbf{Economics example:} \textnormal{After editing \artifact{examples/starter\_article/src/}, hook reminds: ``Re-run \texttt{pytest} and replication-checker before claiming results.''}
\end{frame}

\begin{frame}{Long-Running Goals and Checkpoints}
  \textbf{Problem:} \textnormal{One chat cannot hold a whole pipeline; context drifts.}

  \textbf{Solution:} \textnormal{Explicit goals with \textbf{allowed paths}, \textbf{forbidden paths}, and \textbf{stop conditions}.}
  \vspace{0.3em}
  \begin{tabular}{@{}ll@{}}
    \toprule
    \textbf{Tool} & \textbf{Mechanism} \\
    \midrule
    Codex & \texttt{/goal} or goal UI; \texttt{codex resume} \\
    Claude Code & Checkpoints; \texttt{/goal}; \texttt{claude --continue} (CLI only) \\
    Cursor & Background agents; Agent checkpoints \\
    \bottomrule
  \end{tabular}
  \vspace{1.35em}\\
  
\end{frame}

\begin{frame}{Exercise: Write a Long-Running Goal}
  \textbf{Exercise:} Draft a long-running goal for your project using the following template.
  \vspace{0.8em}

  \begin{block}{Long-Running Goal Template}
    \begin{itemize}
      \item \textbf{Goal:} \textnormal{(Describe your end goal or desired outcome clearly and concretely.)}
      \item \textbf{Allowed paths:} \textnormal{(List productive actions the agent can take to achieve the goal.)}
      \item \textbf{Forbidden paths:} \textnormal{(List actions or areas that are off-limits; e.g., avoid modifying published data, do not email external parties, etc.)}
      \item \textbf{Stop conditions:} \textnormal{(Define clear stopping criteria, e.g., when the goal is achieved, a report is generated, all code passes tests, etc.)}
    \end{itemize}
  \end{block}
  \vspace{1em}
  \textbf{Tip:} Be explicit about boundaries and what success looks like. This helps both you and your agent stay aligned over time.
  \vspace{0.5em}
  {\small Template: \artifact{templates/long\_running\_goal\_prompt.md}}
\end{frame}

\begin{frame}{MCP Servers}
  \textbf{MCP} = Model Context Protocol --- a standard to connect agents to \textbf{external tools}.
  \begin{itemize}
    \item \textbf{What:} \textnormal{Cursor (or other hosts) invoke MCP servers; the agent sees structured tools (e.g.\ FRED API).}
    \item \textbf{Config:} \textnormal{\artifact{examples/starter\_article/.cursor/mcp.json.example} $\to$ \artifact{.cursor/mcp.json}; API keys via environment variables, never in git.}
    \item \textbf{When it helps:} \textnormal{Recurring macro series pulls, scoped filesystem access.}
    \item \textbf{When to skip:} \textnormal{One-off downloads; workshop runs fine without FRED.}
  \end{itemize}
  \vspace{0.25em}
  {\scriptsize\textit{Verify MCP setup in your Cursor version. Docs: \url{https://cursor.com/docs/context/mcp}}}
\end{frame}

\begin{frame}{Skills vs Subagents vs Main Agent}
  \begin{center}
    \begin{tabular}{@{}p{0.28\linewidth}p{0.62\linewidth}@{}}
      \toprule
      \textbf{Use} & \textbf{When} \\
      \midrule
      \textbf{Skill} & \textnormal{Single-purpose procedure (replication-check, referee pass, SDD plan)} \\
      \textbf{Subagent} & \textnormal{Isolated review role; parallel second opinion on memo or PR} \\
      \textbf{Main agent} & \textnormal{Orchestration, merging edits, deciding scope} \\
      \bottomrule
    \end{tabular}
  \end{center}
  \vspace{0.4em}
  \textbf{Invoke explicitly:} \textnormal{``Run replication-checker on repo root'' beats ``check if this replicates.''}
\end{frame}

% =============================================================================
\section{Orchestration}
% =============================================================================

\begin{frame}{Swarms: GitHub as Control Plane}
  A \textbf{swarm} coordinates several agents toward one project outcome --- not one giant prompt.
  \vspace{0.3em}
  \begin{itemize}
    \item \textbf{Issues} = units of work with acceptance criteria and allowed files.
    \item \textbf{Labels:} \artifact{parallel}, \artifact{sequential}, \artifact{blocked-by:\#N}, \artifact{ready-for-review}.
    \item \textbf{Rule:} independent tasks $\to$ parallel agents; dependent tasks $\to$ sequential launches after prior PR merges.
    \item \textbf{Planner} opens issues; implementers work branches; reviewers stay read-only.
  \end{itemize}
  \vspace{0.3em}
  {\small Workshop plan: \artifact{orchestration/card\_krueger\_swarm.md}; issue template: \artifact{orchestration/cloud\_agent\_issue\_template.md}}
\end{frame}

\begin{frame}{Example: Card-Krueger Swarm Streams}
  \footnotesize
  \begin{tabular}{@{}p{0.24\linewidth}p{0.18\linewidth}p{0.22\linewidth}p{0.26\linewidth}@{}}
    \toprule
    \textbf{Stream} & \textbf{Label} & \textbf{Depends on} & \textbf{Reviewer} \\
    \midrule
    Data map audit & sequential start & none & \artifact{data-reviewer} \\
    Cleaning and validation & sequential & data map & \artifact{replication-verifier} \\
    Baseline estimate & sequential & cleaning & \artifact{replication-verifier} \\
    Source note & parallel & source URLs only & \artifact{identification-reviewer} \\
    Teaching write-up & sequential & baseline & \artifact{pr-reviewer} \\
    \bottomrule
  \end{tabular}
  \vspace{0.5em}
  \begin{itemize}
    \item Record handoffs in \artifact{notes/orchestration\_log\_template.md}.
    \item Merge only after review; keep \artifact{main} stable.
    \item Neither swarms nor loops replace human merge approval.
  \end{itemize}
\end{frame}

\begin{frame}{Loop on Verification}
  Three phases, repeated until exit or escalation:
  \begin{enumerate}
    \item \textbf{Implement} --- narrowly scoped change against acceptance criteria.
    \item \textbf{Evaluate} --- \artifact{loop-verifier} subagent or replication-checker skill.
    \item \textbf{Revise} --- address only listed blockers; record iteration in the orchestration log.
  \end{enumerate}
  \vspace{0.3em}
  \textbf{Verdicts:} GREEN $\to$ human review; YELLOW $\to$ loop again; RED or 3 iterations $\to$ stop and open a new issue.
  \vspace{0.3em}
  {\small Protocol: \artifact{orchestration/verification\_loop.md}; skill: \artifact{examples/starter\_article/.cursor/skills/loop-on-verification/}; evaluator: \artifact{examples/starter\_article/.cursor/agents/loop-verifier.md}}
\end{frame}

\begin{frame}{Loop vs Swarm}
  \begin{columns}[T,onlytextwidth]
    \begin{column}{0.48\textwidth}
      \textbf{Swarm}
      \begin{itemize}
        \item Many issues in parallel waves
        \item GitHub labels and dependencies
        \item Merge after subagent review
        \item Example: \artifact{orchestration/card\_krueger\_swarm.md}
      \end{itemize}
    \end{column}
    \begin{column}{0.48\textwidth}
      \textbf{Verification loop}
      \begin{itemize}
        \item One issue, one agent stream
        \item Sequential implement--evaluate cycles
        \item Exit when criteria are green
        \item Skill: \artifact{loop-on-verification}
      \end{itemize}
    \end{column}
  \end{columns}
  \vspace{0.6em}
  Both require acceptance criteria first. Neither replaces human merge approval.
\end{frame}

\begin{frame}{Goal Files and the \texttt{/goal} Command}
  A \textbf{goal} is persistent task state the agent retrieves across sessions --- the local equivalent of a GitHub issue.
  \vspace{0.3em}
  \footnotesize
  \begin{tabular}{@{}p{0.28\linewidth}p{0.66\linewidth}@{}}
    \toprule
    \textbf{Field} & \textbf{Write verifiable criteria} \\
    \midrule
    name & one-line task identifier \\
    description & economic output when done \\
    approach & files to edit, method to use \\
    acceptance\_criteria & checkbox list; each item checkable by command \\
    constraints & raw data untouched; synthetic caveat required \\
    \bottomrule
  \end{tabular}
  \vspace{0.4em}
  {\small Worked examples: \artifact{orchestration/goals/}; template: \artifact{templates/long\_running\_goal\_prompt.md}. Attach via Claude Code \texttt{/goal}, Codex Goals UI, or Cursor Cloud Agent task description.}
\end{frame}

\begin{frame}[t,shrink=8]{Harness Components (Summary)}
  \footnotesize
  \begin{tabular}{@{}p{0.22\linewidth}p{0.72\linewidth}@{}}
    \toprule
    \textbf{Component} & \textbf{What it does (Cursor paths)} \\
    \midrule
    Project context & \textnormal{Always-on: \artifact{AGENTS.md}, \artifact{templates/AGENTS\_economics.md}, \artifact{examples/starter\_article/.cursor/rules/*.mdc}} \\
    Research spec & \textnormal{Plan before code: \artifact{templates/project\_brief.md}, \artifact{templates/research\_design\_memo.md}, \artifact{examples/starter\_article/spec/intent.md}} \\
    Skills & \textnormal{Procedures: \artifact{examples/starter\_article/.cursor/skills/referee-checklist/}, etc.} \\
    Subagents & \textnormal{Isolated reviewers: \artifact{examples/starter\_article/.cursor/agents/}} \\
    Hooks & \textnormal{Deterministic guardrails: \artifact{examples/hooks/.codex/}, \artifact{examples/starter\_article/.cursor/hooks.json}} \\
    MCP & \textnormal{Optional external APIs: \artifact{examples/starter\_article/.cursor/mcp.json.example}} \\
    Gates & \textnormal{Evidence: \artifact{examples/starter\_article/replication/README.md}, \artifact{templates/referee\_report.md}, \artifact{templates/review\_log.md}} \\
    Orchestration log & \textnormal{Lightweight record: what you delegated vs reviewed (\artifact{notes/orchestration\_log\_template.md})} \\
    \bottomrule
  \end{tabular}
\end{frame}

% =============================================================================
\section{Autonomous Agents}
% =============================================================================

\begin{frame}{Autonomous Agents: OpenClaw, Hermes, Eve}
  \begin{tabular}{p{0.21\textwidth}p{0.32\textwidth}p{0.36\textwidth}}
    \toprule
    \textbf{Agent} & \textbf{What to notice} & \textbf{Workshop stance} \\
    \midrule
    OpenClaw & Local-first assistant; channels, tools, skills, sandboxing, multi-agent routing & Good for personal-agent and channel-safety discussion; optional reading \\
    Hermes & Nous Research agent focused on long-running context and user-specific capability growth & Useful for memory, persistence, and ownership of research context \\
    Eve & Vercel framework for building agents with instructions, skills, tools, sandboxes, channels, subagents, schedules & Use only after checking current docs and access boundaries \\
    \bottomrule
  \end{tabular}
  \vspace{0.6em}
  {\small\textbf{Rule:} More autonomy means more need for logs, review gates, and reproducibility evidence.}
\end{frame}

\begin{frame}{Autonomous-Agent Risk Card}
  Fill one risk card before letting any autonomous agent touch a research repo:
  \begin{itemize}
    \item \textbf{Read boundary:} what files, channels, or data can it inspect?
    \item \textbf{Write boundary:} what can it edit, trigger, schedule, or publish?
    \item \textbf{Memory:} what state persists across sessions, and who controls it?
    \item \textbf{Execution:} local, cloud, sandbox, or production system?
    \item \textbf{Review gate:} what human approval is required before outputs enter the repo?
    \item \textbf{Evidence:} what logs, diffs, commands, or replication outputs prove what happened?
  \end{itemize}
  \vspace{0.4em}
  \textbf{Artifact:} \artifact{agent-harness/autonomous\_agent\_risk\_card.md} {\small (see \ParentCourseUrl)}
\end{frame}

\begin{frame}{Risk Card Decision}
  \begin{columns}[T,onlytextwidth]
    \begin{column}{0.48\textwidth}
      \textbf{Approve only if you can answer:}
      \begin{itemize}
        \item \textbf{Read}: \textnormal{what can it inspect?}
        \item \textbf{Write}: \textnormal{what can it edit, trigger, publish, or schedule?}
        \item \textbf{Memory}: \textnormal{what persists and who controls it?}
        \item \textbf{Execution}: \textnormal{local, sandbox, cloud, production, or mixed?}
        \item \textbf{Evidence}: \textnormal{which logs, diffs, transcripts, or outputs prove what happened?}
      \end{itemize}
    \end{column}
    \begin{column}{0.48\textwidth}
      \textbf{Decision labels}
      \begin{itemize}
        \item \textnormal{Safe for read-only use.}
        \item \textnormal{Safe for branch-isolated write use.}
        \item \textnormal{Not safe for this workshop repo.}
        \item \textnormal{Needs more documentation before use.}
      \end{itemize}
      \vspace{0.3em}
      \textbf{Rule}: \textnormal{do not connect private data unless every boundary and approval gate is acceptable.}
    \end{column}
  \end{columns}
\end{frame}

\begin{frame}{Map New Agents Back to the Harness}
  \begin{center}
    \small
    \begin{tabular}{ll}
      \toprule
      \textbf{Autonomous-agent feature} & \textbf{Course harness equivalent} \\
      \midrule
      Markdown instructions & \artifact{AGENTS.md}, \artifact{CLAUDE.md}, Cursor rules \\
      Skills / tools & \artifact{SKILL.md}, MCP tools, validation scripts \\
      Postconditions & Hooks, CI, verification commands \\
      Task memory & Goal files, GitHub issues, orchestration log \\
      Channels / schedules & GitHub issues, CI, scheduled checks \\
      Subagents / workers & data-reviewer, literature-reviewer, loop-verifier \\
      Sandbox & Branch, worktree, remote runner, protected main \\
      Logs & PR description, orchestration log, replication README \\
      \bottomrule
    \end{tabular}
  \end{center}
  \vspace{0.45em}
  \textbf{Teaching point:} The architecture transfers; the UI changes.
\end{frame}

\end{document}
